\documentclass[10pt,a4paper]{amsart}
\usepackage{amsmath,amssymb,amsthm,graphicx}
\usepackage{verbatim}
\newenvironment{reference}{\comment}{\endcomment}
%\newenvironment{reference}{}{}
\newenvironment{slogan}{\comment}{\endcomment}
\newenvironment{history}{\comment}{\endcomment}

% For commutative diagrams we use Xy-pic
\usepackage[all]{xy}

% We use 2cell for 2-commutative diagrams.
\xyoption{2cell}
% Theorem environments.
%
\theoremstyle{plain}
\newtheorem{theorem}[subsection]{Theorem}
\newtheorem{proposition}[subsection]{Proposition}
\newtheorem{lemma}[subsection]{Lemma}

\theoremstyle{definition}
\newtheorem{definition}[subsection]{Definition}
\newtheorem{example}[subsection]{Example}
\newtheorem{exercise}[subsection]{Exercise}
\newtheorem{situation}[subsection]{Situation}

\theoremstyle{remark}
\newtheorem{remark}[subsection]{Remark}
\newtheorem{remarks}[subsection]{Remarks}

\numberwithin{equation}{subsection}

% Macros
%
\def\lim{\mathop{\mathrm{lim}}\nolimits}
\def\colim{\mathop{\mathrm{colim}}\nolimits}
\def\Spec{\mathop{\mathrm{Spec}}}
\def\Hom{\mathop{\mathrm{Hom}}\nolimits}
\def\Ext{\mathop{\mathrm{Ext}}\nolimits}
\def\SheafHom{\mathop{\mathcal{H}\!\mathit{om}}\nolimits}
\def\SheafExt{\mathop{\mathcal{E}\!\mathit{xt}}\nolimits}
\def\Sch{\mathit{Sch}}
\def\Mor{\mathop{\mathrm{Mor}}\nolimits}
\def\Ob{\mathop{\mathrm{Ob}}\nolimits}
\def\Sh{\mathop{\mathit{Sh}}\nolimits}
\def\NL{\mathop{N\!L}\nolimits}
\def\CH{\mathop{\mathrm{CH}}\nolimits}
\def\proetale{{pro\text{-}\acute{e}tale}}
\def\etale{{\acute{e}tale}}
\def\QCoh{\mathit{QCoh}}
\def\Ker{\mathop{\mathrm{Ker}}}
\def\Im{\mathop{\mathrm{Im}}}
\def\Coker{\mathop{\mathrm{Coker}}}
\def\Coim{\mathop{\mathrm{Coim}}}

% Boxtimes
%
\DeclareMathSymbol{\boxtimes}{\mathbin}{AMSa}{"02}

%
% Macros for moduli stacks/spaces
%
\def\QCohstack{\mathcal{QC}\!\mathit{oh}}
\def\Cohstack{\mathcal{C}\!\mathit{oh}}
\def\Spacesstack{\mathcal{S}\!\mathit{paces}}
\def\Quotfunctor{\mathrm{Quot}}
\def\Hilbfunctor{\mathrm{Hilb}}
\def\Curvesstack{\mathcal{C}\!\mathit{urves}}
\def\Polarizedstack{\mathcal{P}\!\mathit{olarized}}
\def\Complexesstack{\mathcal{C}\!\mathit{omplexes}}
% \Pic is the operator that assigns to X its picard group, usage \Pic(X)
% \Picardstack_{X/B} denotes the Picard stack of X over B
% \Picardfunctor_{X/B} denotes the Picard functor of X over B
\def\Pic{\mathop{\mathrm{Pic}}\nolimits}
\def\Picardstack{\mathcal{P}\!\mathit{ic}}
\def\Picardfunctor{\mathrm{Pic}}
\def\Deformationcategory{\mathcal{D}\!\mathit{ef}}
\usepackage{cleveref}
\crefname{originallemma}{original lemma}{original lemmas}
\Crefname{originallemma}{Original lemma}{Original lemmas}
\makeatletter
\newlabel{schemes-lemma-characterize-quasi-separated}{{01KO}{source}{}{sourcepointer-1}{}}
\newlabel{schemes-lemma-characterize-quasi-separated@cref}{{[originallemma][0][]01KO}{[1][0][]source}{}{}{}}
\makeatother
\crefname{originallemma}{original lemma}{original lemmas}
\Crefname{originallemma}{Original lemma}{Original lemmas}
\makeatletter
\newlabel{schemes-lemma-valuative-criterion-separatedness}{{01L0}{source}{}{sourcepointer-2}{}}
\newlabel{schemes-lemma-valuative-criterion-separatedness@cref}{{[originallemma][0][]01L0}{[1][0][]source}{}{}{}}
\makeatother
\title{A quasi-compact quasi-separated scheme with no separated quasi-compact dense open}
\author{The Stacks Project Authors}
\date{stacks-project:examples.tex:local-snapshot-20261012; GNU FDL 1.2 or later}
\begin{document}
\maketitle
\noindent\textit{Editorial note.} All delivered mathematical text is the author's own, copied from the pinned Stacks Project chapter file examples.tex; no statement or proof step was rewritten, reordered or removed.\par\medskip
\noindent\textit{Editorial note.} The counted proof is the author's complete argument: The argument's final contradiction is the author's own 'by the first paragraph of this section'; the first paragraph is included as the first context span.\par\medskip
\noindent\textit{Editorial note.} The label lemma-no-dense-separated-quasi-compact-open-in-qcqs names the counted statement, and the proof environment immediately following it in the source is delivered verbatim as its proof.\par\medskip
\noindent\textit{Editorial note.} The counted statement is carried by the author's own lemma environment in the chapter source.\par\medskip
\noindent\textit{Allowed standard prerequisites (editorial).}\begin{itemize}
\item Context span 2249-2274 (delivered as L48311): no quasi-compact open dense in S - \{s\}, and the T/U complement statement, including the isomorphism (T, t) = (S, s) sending x to z\_1 and z\_i to z\_\{i+1\}
\item Context span 2293-2312: X is the glueing of two copies of T along the affine open U; pi: X -> T maps U\_i isomorphically to T and U\_1 cap U\_2 isomorphically to U; X is quasi-separated (Schemes, characterize-quasi-separated) and quasi-compact; for the closed points x\_1 in U\_1, x\_2 in U\_2 over t and the generic point etatilde in U\_1 cap U\_2 over eta, we have etatilde leads to x\_1 and etatilde leads to x\_2; since pi|\_W: W -> T is separated, W cannot contain both x\_1 and x\_2 (valuative criterion, Schemes, valuative-criterion-separatedness); say x\_1 not in W; then W cap U\_1 is a quasi-compact (X quasi-separated) dense open of U\_1 not containing x\_1; via (T, t) = (S, s) this contradicts the first context span
\end{itemize}\medskip
\ifdefined\hypertarget\hypertarget{sourcepointer-1}{}\fi
\noindent\textit{Source reference.} Original lemma 01KO: sources/stacks-chow/tags.txt, source lines 2057--2057. Prerequisite pointer invoked at 2353; it names a standard fact of a sibling chapter that the counted argument uses, and every step of the counted proof built on it is retained verbatim in the delivered body.\par\medskip
\ifdefined\hypertarget\hypertarget{sourcepointer-2}{}\fi
\noindent\textit{Source reference.} Original lemma 01L0: sources/stacks-chow/tags.txt, source lines 2069--2069. Prerequisite pointer invoked at 2363; it names a standard fact of a sibling chapter that the counted argument uses, and every step of the counted proof built on it is retained verbatim in the delivered body.\par\medskip
Let $k$ be a field and let $A = k[z_1, z_2, z_3, \ldots]/I$ where $I$
is the ideal generated by all pairwise products
$z_iz_j$, $i \not = j$, $i, j \in \mathbf{N}$. Set $S = \Spec(A)$.
Let $s \in S$ be the closed point corresponding to the maximal
ideal $(z_i)$. We claim there is no
quasi-compact open $V \subset S \setminus \{s\}$ which is
dense in $S \setminus \{s\}$. Note that $S \setminus \{s\} = \bigcup D(z_i)$.
Each $D(z_i)$ is open and irreducible with generic point
$\eta_i$. We conclude that $\eta_i \in V$ for all $i$.
However, a principal affine open of $S \setminus \{s\}$ is
of the form $D(f)$ where $f \in (z_1, z_2, \ldots)$. Then
$f \in (z_1, \ldots, z_n)$ for some $n$ and we see that $D(f)$ contains
only finitely many of the points $\eta_i$. Thus $V$ cannot be quasi-compact.

\medskip\noindent
Let $k$ be a field and let $B = k[x, z_1, z_2, z_3, \ldots]/J$ where $J$
is the ideal generated by the products $xz_i$, $i \in \mathbf{N}$ and by
all pairwise products $z_iz_j$, $i \not = j$, $i, j \in \mathbf{N}$.
Set $T = \Spec(B)$. Consider the principal open $U = D(x)$.
We claim there is no quasi-compact open $V \subset T$ such that
$V \cap U = \emptyset$ and $V \cup U$ is dense in $T$.
Let $t \in T$ be the closed point corresponding to the maximal
ideal $(x, z_i)$. The closure of $U$ in $T$ is
$\overline{U} = U \cup \{t\}$. Hence $V \subset \bigcup_i D(z_i)$
is a quasi-compact open. By the arguments of the previous paragraph
we see that $V$ cannot be dense in $\bigcup D(z_i)$.

Let $X$ be the glueing of two copies of the affine scheme $T$ (see above)
along the affine open $U$. Thus there is a morphism $\pi : X \to T$ and
$X = U_1 \cup U_2$ such that $\pi$ maps $U_i$ isomorphically to $T$ and
$U_1 \cap U_2$ isomorphically to $U$. Note that $X$ is quasi-separated
(by Schemes, Lemma \ref{schemes-lemma-characterize-quasi-separated})
and quasi-compact. We claim there does not exist a separated, dense,
quasi-compact open $W \subset X$. Namely, consider the two closed
points $x_1 \in U_1$, $x_2 \in U_2$ mapping to the closed point $t \in T$
introduced above. Let $\tilde \eta \in U_1 \cap U_2$ be the generic point
mapping to the (unique) generic point $\eta$ of $U$.
Note that $\tilde\eta \leadsto x_1$ and $\tilde\eta \leadsto x_2$
lying over the specialization $\eta \leadsto t$.
Since $\pi|_W : W \to T$ is separated we conclude that we cannot have both
$x_1$ and $x_2 \in W$ (by the valuative criterion of separatedness
Schemes, Lemma \ref{schemes-lemma-valuative-criterion-separatedness}).
Say $x_1 \not \in W$. Then $W \cap U_1$ is a quasi-compact (as $X$
is quasi-separated) dense open of $U_1$ which does not contain $x_1$.
Now observe that there exists an isomorphism $(T, t) \cong (S, s)$
of schemes (by sending $x$ to $z_1$ and $z_i$ to $z_{i + 1}$).
Hence by the first paragraph of this section we arrive at a contradiction.

\begin{lemma}
\label{lemma-no-dense-separated-quasi-compact-open-in-qcqs}
There exists a quasi-compact and quasi-separated scheme $X$ which does
not contain a separated quasi-compact dense open.
\end{lemma}

\begin{proof}
See discussion above.
\end{proof}

\end{document}
